\section{System 5: The Resonant Interface (Electroweak Scale) \texorpdfstring{($v, \lambda, m_H, y_e$)}{(v, lambda, mH, ye)}}
\label{sec:resonant_interface}

\CatchFileBetweenTags{\AlphaInvVal}{constants.tex}{AlphaInvVal}
\CatchFileBetweenTags{\CompMAR}{constants.tex}{CompMAR}
\CatchFileBetweenTags{\MeMeVPrint}{constants.tex}{MeMeVPrint}

% CAP
\CatchFileBetweenTags{\HiggsVEVStepOneVal}{constants.tex}{HiggsVEVStepOneVal}
\CatchFileBetweenTags{\HiggsVEVStepTwoVal}{constants.tex}{HiggsVEVStepTwoVal}
\CatchFileBetweenTags{\HiggsVEVVal}{constants.tex}{HiggsVEVVal}
\CatchFileBetweenTags{\HiggsVEVExperimentalValue}{constants.tex}{HiggsVEVExperimentalValue}
\CatchFileBetweenTags{\HiggsVEVAccText}{constants.tex}{HiggsVEVAccText}
\CatchFileBetweenTags{\HiggsImpedanceVal}{constants.tex}{HiggsImpedanceVal}
\CatchFileBetweenTags{\WeakApertureProjVal}{constants.tex}{WeakApertureProjVal}

\CatchFileBetweenTags{\FermiConstVal}{constants.tex}{FermiConstVal}
\CatchFileBetweenTags{\FermiConstExperimentalValue}{constants.tex}{FermiConstExperimentalValue}
\CatchFileBetweenTags{\FermiConstAccText}{constants.tex}{FermiConstAccText}

% PRO
\CatchFileBetweenTags{\HiggsLambdaVal}{constants.tex}{HiggsLambdaVal}
\CatchFileBetweenTags{\HiggsLambdaExperimentalValue}{constants.tex}{HiggsLambdaExperimentalValue}
\CatchFileBetweenTags{\HiggsLambdaAccText}{constants.tex}{HiggsLambdaAccText}

% TOL
\CatchFileBetweenTags{\HiggsMassVal}{constants.tex}{HiggsMassVal}
\CatchFileBetweenTags{\HiggsMassExperimentalValue}{constants.tex}{HiggsMassExperimentalValue}
\CatchFileBetweenTags{\HiggsMassAccText}{constants.tex}{HiggsMassAccText}

% MAR
\CatchFileBetweenTags{\ElectronYukawaVal}{constants.tex}{ElectronYukawaVal}
\CatchFileBetweenTags{\ElectronYukawaStepOneVal}{constants.tex}{ElectronYukawaStepOneVal}
\CatchFileBetweenTags{\ElectronYukawaExperimentalValue}{constants.tex}{ElectronYukawaExperimentalValue}
\CatchFileBetweenTags{\ElectronYukawaAccText}{constants.tex}{ElectronYukawaAccText}

\subsection{The Architectural Requirement}
The vacuum must instantiate a nested persistent subsystem to translate information between the high-frequency temporal lattice resonance ($\Delta=43$) and the specific topological boundary of the Weak Interaction ($\chi=2$) without reflection loss. In the Standard Model, this manifests as the Higgs sector ($v, \lambda, \mu^2, y_e$), whose specific values constitute the Hierarchy Problem. Here, we derive them not as arbitrary tuning parameters, but as the exact numerical capacity allocations of a nested impedance-matching subsystem.

As derived in \textit{Informational Energetics: Entropic Action} \cite{meyer_ieaction_2026}, the architecture of persistence mathematically necessitates the Higgs functional form to bound the finite state-space capacity of the vacuum. In this section, we assume those derived functional forms and focus exclusively on calculating their exact numerical partition coefficients using the nested partitioning rules. 

To satisfy \textbf{Rule 3: Impedance Matching (Shannon's Channel Theorem)}, the internal bandwidth of this resonant interface is rigidly partitioned across the six pillars using the lattice invariants:

\begin{itemize}
    \item \textbf{Capacity -- The VEV ($v$):} The mass baseline allocated to the resonant interface, against which the vacuum potential is normalized.
    \item \textbf{Map -- Scalar Hypercharge ($Y=1$):} The unitary address. It corresponds to the inverse of the lattice Unitary Ground State ($\Delta^0=1$).
    \item \textbf{Protocol -- Self-Coupling ($\lambda$):} The self-interaction bandwidth allocation, uniquely determined by \textbf{Rule 1: Capacity Partitioning}.
    \item \textbf{Governor -- The Quartic Potential ($\lambda|\phi|^4$):} The geometric capacity wall enforcing the topological boundary ($\chi=2$).
    \item \textbf{Toll -- Instability Factor ($\mu^2$):} The fractional entropic cost of maintaining the broken symmetry state.
    \item \textbf{Margin -- The Electron Yukawa ($y_e$):} The minimal resolvable noise floor, structurally determined by \textbf{Rule 2: Reversible Encoding}.
\end{itemize}

\subsection{Capacity: The Higgs VEV (\texorpdfstring{$v$}{v})}
We calculate the static potential minimum strictly allocated by the lattice invariants.

\subsubsection{The Bare Geometric Floor}
The bare VEV is derived from the total causal area of the topological boundary ($\chi \Delta^2$) minus the structural overhead of the substrate ($L_{substrate}$). 
\begin{equation}
v_{geo} = (\chi \Delta^2 - L_{substrate}) \cdot \alpha^{-1} \cdot m_e 
\end{equation}

Dimensionally, the geometric terms $(\chi \Delta^2 - L_{substrate})$ and $\alpha^{-1}$ strictly count discrete state capacities and transmission ratios (dimensionless scalars). To map this raw information capacity into a physical energy scale, it must be multiplied by the fundamental mass bit $m_e$ (the geometric resolution floor of the lattice, establishing the [Mass] dimension).

Substituting the Substrate Load ($L_{substrate} = 188$, defined in Paper 1 \cite{meyer_informational_2026} as the total information the vacuum must process per fundamental cycle: $\nu + (D \cdot \Delta)$):
\begin{equation}
\begin{split}
v_{geo} &= (2 \cdot 43^2 - 188) \cdot \AlphaInvVal \cdot \MeMeVPrint \text{ MeV} \\
&\approx \HiggsVEVStepOneVal
\end{split}
\end{equation}

\subsubsection{The Geometric Corrections}
The physical vacuum is not an abstract background; it is a topological manifold subject to thermodynamic fluctuations. To bridge the ideal geometric floor ($v_{geo}$) to physical reality, we apply two independent scale factors mandated by the lattice projection. Because these represent fractional boundary and thermal effects, they manifest as exact leading-order linear expansions $(1 \pm x)$. Higher-order geometric deviations are strictly suppressed by $\alpha^2 \approx 5 \times 10^{-5}$, safely truncating the expansion.

\paragraph{Topological Dilution (Spatial Geometry)}
The regulator energy is diluted by the topology of the charge boundaries. We define the \textbf{Effective Coupling Volume} ($D_{eff}$) of the space as the manifold rank ($D=4$) augmented by the topological boundary ($\chi=2$) distributed over the spherical gauge geometry ($4\pi$).
\begin{equation}
D_{eff} = D + \frac{\chi}{4\pi} \approx 4.15915
\end{equation}
\textit{(Scope Defense: Adding $\chi/4\pi$ does not break the integer dimensionality of the manifold; rather, the continuous $4\pi$ geometric factor represents the exact macroscopic continuum limit of the solid angle of the $\chi=2$ boundary, acting as a continuous geometric correction to the coupling volume).}

This creates a first-order geometric dilution factor based on the vacuum coupling:
\begin{equation}
\delta_{topo} = \left( 1 + \frac{\alpha}{D_{eff}} \right)
\end{equation}

\paragraph{Thermodynamic Noise (Temporal Stability)}
We must also account for the lattice's thermal noise floor. As derived formally in Section \ref{subsec:margin_electron}, the \textbf{Margin Impedance} ($Z_{MAR} \approx \CompMAR$) represents the absolute baseline resolution of the substrate. Since the vacuum stability floor supports exactly 3 fermion generations (the available residual degrees of freedom: $\sigma - \chi = 5 - 2 = 3$), this noise amplitude is partitioned linearly across these 3 slots. This creates the leading-order thermodynamic reduction factor:
\begin{equation}
\delta_{noise} = \left( 1 - \frac{Z_{MAR}}{\sigma - \chi} \right)
\end{equation}

\subsubsection{Synthesis: The Physical VEV}
The physical Vacuum Expectation Value is the geometric floor modulated by the simultaneous application of these topological and thermodynamic factors.
\begin{equation}
v_{phys} = v_{geo} \cdot \delta_{topo} \cdot \delta_{noise}
\end{equation}

Substituting the derived values:
\begin{equation}
\begin{split}
v_{phys} &= (\HiggsVEVStepOneVal) \cdot \left( 1 + \frac{\alpha}{4.159} \right) \cdot \left( 1 - \frac{\CompMAR}{3} \right) \\
& \approx \mathbf{\HiggsVEVVal}
\end{split}
\end{equation}

\begin{itemize}
    \item \textbf{Experimental Value:} \HiggsVEVExperimentalValue
    \item \textbf{Accuracy:} \HiggsVEVAccText
\end{itemize}

\subsubsection{Derived Limit: The Fermi Constant (\texorpdfstring{$G_F$}{GF})}
In the Standard Model, the effective 4-fermion weak interaction cross-section contains a $1/\sqrt{2}$ normalization factor. In standard QFT, this is treated as an algebraic convention. In the $E_8$-Persistence theory, it is a strict geometric projection.

The weak interaction mediates flux across the 2D spherical topological boundary ($\chi=2$) of the node. For this 2D isotropic boundary flux to interact with a 1D chiral fermion trajectory, the flux must be projected onto the 1D interaction axis. The geometric projection of a 2D isotropic amplitude onto a 1D interaction axis requires taking the root-mean-square (RMS) projection, yielding exactly $1/\sqrt{2}$.

Therefore, the Fermi constant ($G_F$) is precisely the inverse squared VEV attenuated by this exact topological flux projection:
\begin{equation}
    G_F = \frac{1}{\sqrt{2}v_{phys}^2} \approx \mathbf{\FermiConstVal}
\end{equation}

\begin{itemize}
\item \textbf{Experimental Value:} \FermiConstExperimentalValue
\item \textbf{Accuracy:} \FermiConstAccText
\end{itemize}

\subsection{Identity (\texorpdfstring{$\Delta^0$}{Delta0}): Scalar Hypercharge (\texorpdfstring{$Y=1$}{Y1})} 
The address of the regulator corresponds to the inverse of the lattice Unitary Ground State ($\Delta^0 =1$). This unitary charge, combined with the topological boundary ($\chi=2$), geometrically mandates an $SU(2)$ doublet structure to interface with the chiral matter sector. (\textit{Note: A complete derivation of how all Standard Model fermion hypercharges emerge as exact geometric ratios of the lattice invariants is provided in Appendix \ref{app:OriginOfHypercharge}}).

\subsection{Protocol: Self-Coupling \texorpdfstring{$\lambda$}{lambda} (Capacity Partitioning)}
In Informational Energetics, coupling constants are strictly bandwidth allocations. Under \textbf{Rule 1: Capacity Partitioning}, the self-coupling $\lambda$ is defined as the exact ratio of the \textbf{Net Available Bandwidth} to the \textbf{Total Systemic Channel Capacity}.

\subsubsection{The Net Bandwidth}
The bandwidth available for the scalar sector is the Interaction Remainder ($\sigma - \chi$) minus the \textbf{Resonant Cost} ($1/\Delta$).
\begin{equation}
\text{Net} = (\sigma - \chi) - \frac{1}{\Delta} = 3 - \frac{1}{43} \approx 2.9767
\end{equation}
The factor $1/\Delta$ represents the fractional bandwidth consumed by temporal synchronization. Without this reservation, the field would drift out of phase with the lattice update cycle, causing decoherence.

\subsubsection{The Coupling Calculation}
\begin{equation}
\lambda = \frac{\text{Net}}{L_{intrinsic}} = \frac{2.9767}{23} \approx \mathbf{\HiggsLambdaVal}
\end{equation}
\begin{itemize}
    \item \textbf{Experimental Value:} \HiggsLambdaExperimentalValue
    \item \textbf{Accuracy:} \HiggsLambdaAccText
\end{itemize}

\textbf{Physical Interpretation:} The Higgs self-coupling is not arbitrary; it is exactly the specific fraction of vacuum bandwidth ($\approx 13\%$) remaining for self-regulation after paying the entropic overhead for temporal synchronization.

\subsection{Toll: Instability \texorpdfstring{$\mu^2$}{mu2}}
The instability factor $\mu^2$ drives Spontaneous Symmetry Breaking. This is the \textbf{Toll}, the entropic cost of maintaining the False Vacuum state. 

The $SU(2)$ doublet structure ($\chi=2$) requires two independent complex fields. The total instability action ($T_H$) geometrically mandated by the substrate is the sum of their individual capacity allocations ($2 \times \lambda$):
\begin{equation}
T_H = \chi \cdot \lambda = 2\lambda \approx 0.259
\end{equation}
We use this derived geometric Toll to structurally evaluate the physical mass parameter $\mu^2$. Scaling the VEV capacity by the fractional Toll per degree of freedom ($T_H/2$) yields:
\begin{equation}
\mu^2 = \frac{T_H}{2} v^2 = \frac{2\lambda}{2} v^2 = \lambda v^2
\end{equation}
This acts as a structural consistency check, successfully recovering the Standard Model relation ($\mu^2 = \lambda v^2$) from geometric first principles and proving that the negative tachyonic mass parameter $-\mu^2$ is not a mathematical trick, but the strict manifestation of the entropic Toll acting on the VEV.

\subsubsection{Output: The Higgs Mass (\texorpdfstring{$m_H$}{mH})}
With the Capacity ($v$) and Protocol ($\lambda$) established, the mass of the scalar excitation follows strictly from the closure of the geometric system.
\begin{equation}
\begin{split}
m_H &= \sqrt{2\lambda} v_{phys} = \sqrt{2 (\HiggsLambdaVal)} \cdot (\HiggsVEVVal) \\
& \approx \mathbf{\HiggsMassVal}
\end{split}
\end{equation}
\begin{itemize}
    \item \textbf{Experimental Value:} \HiggsMassExperimentalValue
    \item \textbf{Accuracy:} \HiggsMassAccText
\end{itemize}

\subsection{Margin: The Electron Yukawa \texorpdfstring{$y_e$}{ye} (Reversible Encoding)}
\label{subsec:margin_electron}
The Margin represents the smallest non-zero bit of mass the lattice can resolve against thermal noise. Under \textbf{Rule 2: Reversible Encoding (Landauer's Principle)}, the system must conserve mismatch by storing it as structural metadata. We derive this minimal scale as the Unit Bit ($1$) diluted over the \textbf{Total Configuration Space Volume}. 

This establishes the Margin Impedance ($Z_{MAR}$)\cite{meyer_informational_2026} for the substrate:
\begin{equation}
y_{e,bare} \equiv Z_{MAR} = \frac{1}{V_{config}} = \frac{1}{\underbrace{L_{embed}}_{\text{Budget}} \cdot \underbrace{(\sigma+1)}_{\text{Aperture}} \cdot \underbrace{\Delta^2}_{\text{Area}}}
\end{equation}

\subsubsection{Physical Interpretation: The Landauer Limit}
Physically, the electron represents the \textbf{Landauer Limit} of the vacuum geometry. Any topological knot with a coupling $y < y_{e,bare}$ would possess a binding energy lower than the thermal noise floor of the substrate ($E < k_B T_{eff} \ln 2$), making it informationally indistinguishable from vacuum fluctuations. The electron is stable because it is the \textbf{Minimum Viable Particle}—the physical instantiation of a single, non-erasable bit of mass.

\subsubsection{The Geometric Self-Energy Correction}
Standard field theory treats the difference between bare and physical parameters as ``Renormalization'' arising from virtual particle loops. In the $E_8$-Persistence theory, this difference is the \textbf{Geometric Stress} of embedding the knot's internal symmetry into the spacetime manifold.

A charged particle possesses an internal interaction structure of rank $\sigma=5$, but it must exist within a manifold of rank $D=4$. This dimensional mismatch ($\sigma > D$) creates an irreducible ``Overpressure'' on the surrounding lattice. We define the \textbf{Projection Coefficient} ($\mathcal{P}$) as the strict geometric ratio of this internal load to the manifold capacity:
\begin{equation}
\mathcal{P} = \frac{\text{Interaction Rank}}{\text{Manifold Rank}} = \frac{\sigma}{D} = \frac{5}{4} = 1.25
\end{equation}
The physical effective coupling ($y_{phys}$) is the bare geometric floor ($y_{bare}$) enhanced by this projection stress acting through the vacuum impedance ($\alpha$):
\begin{equation}
y_{e,phys} = y_{e,bare} \cdot \left(1 + \mathcal{P} \cdot \alpha \right) = y_{e,bare} \cdot \left(1 + \frac{\sigma}{D}\alpha \right)
\end{equation}
This represents the exact leading-order geometric correction. Higher-order structural feedback terms are suppressed by $\alpha^2 \approx 5 \times 10^{-5}$, strictly grounding the truncation in the same dimensional budget logic utilized throughout the framework.

Substituting the invariants ($L_{embed}=31, \sigma+1=6, \Delta=43$) and the derived $\alpha$:
\begin{equation}
\begin{split}
y_{e,phys} &\approx \left(\frac{1}{31 \cdot 6 \cdot 43^2}\right) \cdot \left(1 + 1.25 \cdot \frac{1}{137.036}\right) \\
&\approx \ElectronYukawaStepOneVal \cdot (1.00912) \\
&\approx \mathbf{\ElectronYukawaVal}
\end{split}
\end{equation}

\begin{itemize}
    \item \textbf{Experimental Value:} \ElectronYukawaExperimentalValue
    \item \textbf{Accuracy:} \ElectronYukawaAccText
\end{itemize}

\subsubsection{Physical Consequences: The Spectral Floor}
This geometric floor establishes the lower bound of information storage in the vacuum, explaining several fundamental features of the particle spectrum:
\begin{itemize}
    \item \textbf{Why the electron is stable:} It sits exactly at the minimum resolvable mass. There is no ``lower shelf'' to decay to.
    \item \textbf{Why neutrinos are so light:} Massive neutrinos possess Yukawa couplings $y_\nu \ll y_e$, pushing them below the vacuum resolution floor. This structurally prohibits them from acquiring mass via the standard Higgs mechanism (which strictly terminates at $y_e$), necessitating the alternative phonon mechanism (derived in subsequent work).
\end{itemize}

\subsection{Validation: Global Impedance Matching}
The ultimate structural test of the Resonant Interface is whether it successfully satisfies \textbf{Rule 3: Impedance Matching}. If the Higgs mechanism is a functional regulator, its total internal impedance ($Z_H$) must perfectly match the \textbf{Projected Aperture} of the force it mediates (the Weak Interaction).

\begin{enumerate}
    \item \textbf{Higgs Impedance ($Z_H$):} Defined as the Inverse Protocol ($1/\lambda$). Crucially, because it operates under the Entropic Action survivorship filter ($P \propto e^{-S_\Phi}$) established in Paper 1, and its action cost is $T_H = 2\lambda$, this impedance is exponentially modulated by its Toll:
    \begin{equation}
    Z_H = \frac{1}{\lambda} e^{-2\lambda} \approx \frac{1}{0.129424} (0.7719) \approx \mathbf{\HiggsImpedanceVal}
    \end{equation}
    
    \item \textbf{Weak Projected Aperture ($A_{weak}$):} The base geometric aperture of the interaction is the sum of the lattice primitives (links and nodes) involved: ($\sigma + 1 = 6$). However, because this projection occurs across the physical manifold, it must pay the \textbf{The Coordinate Overhead} cost ($\eta$):
    \begin{equation}
    A_{weak} = (\sigma + 1) \cdot \eta \approx \mathbf{\WeakApertureProjVal}
    \end{equation}
\end{enumerate}

\textbf{Result:} The regulator impedance (\HiggsImpedanceVal) matches the projected aperture (\WeakApertureProjVal) to within \textbf{0.01\%}. 

This structural lock confirms that the Electroweak Scale is not an arbitrary tuning parameter, but is derived entirely from the lattice resonance satisfying global impedance matching. The $O(10^{-4})$ residual is well within the truncation bounds of the leading-order geometric corrections (suppressed by $\alpha^2 \sim 5 \times 10^{-5}$) and the precision limits of the empirical mass anchors, validating the closure of the Resonant Interface.